




We here provide technical preliminaries.
We introduce extra notation in \Cref{sec:notation}. In \Cref{sec:uhl intro}, we overview known facts about the Uhlmann transformation.
We present the three computational frameworks for quantum algorithm design in \Cref{sec:three computational models}.
In \Cref{sec:previous key subroutine}, we give a quick review of key algorithmic primitives, including block encoding, the QSVT, and the density matrix exponentiation, which form the building blocks of our algorithms.



\subsection{Extra notations}
\label{sec:notation}


We denote a Hilbert space by $\cH$, and the space associated with a system $\sA$ is denoted by $\cH^\sA$. Hilbert spaces such as $\cH^{\sA'}$, $\cH^{\hat{\sA}}$, or $\cH^{\sA_1}$ are isomorphic to $\cH^\sA$, meaning that they have the same dimension and can be equipped with the same orthonormal basis as $\cH^\sA$.
This applies not only to the system $\sA$, but also to any systems, such as $\cH^{\sB'}$ and $\cH^{\hat{\sC}}$. The dimension of $\cH$ is denoted by $d$ and, for instance, the dimension of $\cH^\sA$ is denoted by $d_\sA$. 




For an operator $M$, the kernel of $M$, i.e., the subspace spanned by vectors mapped to zero by $M$, is denoted by $\ker[M]$. The support of $M$ is denoted by $\supp[M]$, which is defined as the orthogonal complement of $\ker[M]$.
The image of $M$, i.e., the subspace spanned by all vectors of the form $M\ket{v}$, is denoted by $\im[M]$.
For any operator $M$, we denote the complex conjugate and the transpose in a given orthonormal basis by $M^*$ and $M^\top$, respectively, and denote the Hermitian conjugate by $M^\dag$. 





We omit the symbol of the tensor product between vectors and denote it as $\ket{\rho}\otimes\ket{\sigma} = \ket{\rho}\ket{\sigma}$.
We denote by $\ket{\Phi}$ the maximally entangled state defined in the orthonormal computational basis. For instance, the maximally entangled state between $\sA$ and $\hat{\sA}$ is $\ket{\Phi}^{\sA\hat{\sA}} = \f{1}{\sqrt{d_\sA}}\sum_{i=1}^{d_\sA}\ket{i}^\sA\ket{i}^{\hat{\sA}}$, where $\{\ket{i}\}_{i=1}^{d_\sA}$ is the computational basis in $\sA$ and $\hat{\sA}$, respectively. 
We also denote an unnormalized maximally entangled state by $\ket{\Gamma}$, e.g., $\ket{\Gamma}^{\sA\hat{\sA}} = \sum_{i=1}^{d_\sA}\ket{i}^\sA\ket{i}^{\hat{\sA}}$ for $\sA$ and $\hat{\sA}$.
Note that $\sum_{i=1}^d\ket{i}\ket{i} = \sum_{i=1}^d\ket{e_i}\ket{e_i^*}$ holds for any orthonormal basis $\{\ket{e_i}\}_{i=1}^d$, where the complex conjugate is taken with respect to the computational basis~$\{\ket{i}\}_{i=1}^d$.


For an isometry $U$, the corresponding isometry channel $\cU$ is defined by $\cU(\cdot) = U(\cdot)U^\dag$, and its inverse is defined by $\cU^\dag = U^\dag(\cdot)U$.
For any state $\psi^\sA$, let $\cP_{\psi}^{\bC\rarr\sA}$ be the state-preparation channel defined by $\cP_{\psi}^{\bC\rarr\sA} = \psi^\sA$.
We denote the trace map by $\tr$, and denote the partial trace map over a subsystem, for example, $\sA$, by $\tr_\sA$.
We denote the swap operator by $F$, such as $F^{\sA\hat{\sA}}$, which swaps the systems $\sA$ and $\hat{\sA}$.







For a matrix $M$, the Schatten-$p$ norm is defined by $\| M \|_p = \big(\tr\big[\big(\sqrt{M^\dag M}\big)^p\big]\big)^{1/p}$ for $p\in[1, \infty]$. We particularly use the trace norm, the Hilbert--Schmidt norm, and the operator norm, corresponding to $p=1, 2, \infty$, respectively.
The Schatten $p$-norm is isometric invariant~\cite{wilde2013QItheory, bhatia2013matrix, watrous2018TheoryQI}: $\|M\|_p = \|UMV\|_p,$ for isometries $U$ and $V$, and is monotonic: $\|M\|_p \leq \|M\|_q$, for $p \geq q$.
For the Schatten $p$-norm, the H\"{o}lder's inequality~\cite{Rogers1888extension, Holder1889inequality} holds: for any matrix $M_1$ and $M_2$, we have that
\begin{align}
\label{eq:Holder ineq}
    \|M_1M_2\|_r \leq \|M_1\|_p\|M_2\|_q,
\end{align}
where $1/r = 1/p + 1/q$.
The trace norm has a contraction property against the partial trace~\cite{rastegin2012NormPartialTrace, watrous2018TheoryQI} such that for any matrix $M^{\sA\sB}$,
\begin{equation}
\label{eq:contraction trace norm}
    \big\|M^{\sA\sB}\big\|_1 \geq \big\|\tr_{\sB}\big[M^{\sA\sB}\big]\big\|_1.
\end{equation}
Trace norm also has a variational characterization~\cite{nielsen2010quantum, wilde2013QItheory, watrous2018TheoryQI}: for any quantum states $\rho$ and $\sigma$, it holds that
\begin{align}
\label{eq:trace distance variational}
    \f{1}{2}\|\rho - \sigma\|_1 = \max_{M; 0 \leq M \leq \bI}\tr[M(\rho - \sigma)].
\end{align}
For the trace norm and the Hilbert--Schmidt norm, the Powers--St\o rmer inequality~\cite{powers1970free, kittaneh1987inequalities, bhatia2013matrix} holds: for any Hermitian matrices $M_1$ and $M_2$, 
\begin{align}
\label{eq:powers stormer ineq}
    \big\| M_1^{1/2} - M_2^{1/2} \big\|_2 \leq \big\|M_1 - M_2 \big\|_1^{1/2}.
\end{align}




The diamond norm~\cite{Kitaev1997QCAlgoEC, watrous2004notessuperoperatornorm, watrous2018TheoryQI} for a linear map $\cT^{\sA\rarr\sB}$ is defined by
\begin{align}
    \big\|\cT^{\sA\rarr\sB}\big\|_\diamond = \max_{\sR}\max_{M^{\sR\sA}; \|M^{\sR\sA}\|_1 \leq 1}\big\|\rid^\sR \otimes \cT^{\sA\rarr\sB}(M^{\sR\sA})\big\|_1, 
\end{align}
where the maximization is taken over all reference systems $\sR$ of arbitrary size and all operators $M^{\sR\sA}$ with trace norm at most one.
For quantum channels, it suffices to take the maximization over all pure states on $\sR\sA$, where $\sR$ can be chosen such that $d_\sR = d_\sA$.
The diamond norm is subadditive for the composition of quantum channels~\cite{watrous2018TheoryQI}: for any quantum channels $\cT_1^{\sA\rarr\sB}$, $\cT_2^{\sB\rarr\sC}$, $\cE_1^{\sA\rarr\sB}$, and $\cE_2^{\sB\rarr\sC}$, it holds that
\begin{align}
    &\big\|\cT_2^{\sB\rarr\sC} \circ \cT_1^{\sA\rarr\sB} - \cE_2^{\sB\rarr\sC} \circ \cE_1^{\sA\rarr\sB}\big\|_\diamond     \leq \big\|\cT_2^{\sB\rarr\sC} - \cE_2^{\sB\rarr\sC}\big\|_\diamond + \big\|\cT_1^{\sA\rarr\sB} - \cE_1^{\sA\rarr\sB}\big\|_\diamond.
\end{align}





The trace norm and the fidelity are related by the Fuchs--van de Graaf inequalities~\cite{Fuchs1999fuchsvandegraaf}: 
\begin{equation}
\label{eq:fuchs van de graaf}
   1-\sqrt{\rF}(\rho, \sigma) \leq \f{1}{2}\|\rho - \sigma \|_1 \leq \sqrt{1-\rF(\rho, \sigma)}. 
\end{equation}
For isometry channels $\til{\cV}$ and $\cV$, the diamond norm and the operator norm are related as~\cite{Kitaev1997QCAlgoEC} 
\begin{align}
\label{eq:rel of operator and diamond norm}
    \f{1}{2}\big\|\til{\cV} - \cV\big\|_\diamond \leq \big\|\til{V} - V\big\|_\infty.
\end{align}
The diamond norm and the fidelity satisfy the following relation: for any state $\rho$ and any pure state $\ket{\phi}$,
\begin{align}
\label{eq:fidelity and diamond}
    \big|\rF\big(\cT_1(\rho), \ket{\phi}\big) - \rF\big(\cT_2(\rho), \ket{\phi}\big)\big| \leq \f{1}{2}\big\|\cT_1 - \cT_2\big\|_\diamond,
\end{align}
which follows from the variational characterization of the trace norm and the definition of the diamond norm.


For a vector $v \in \bC^d$, the Euclidean norm is defined by $\|v\| =\sqrt{v^\dag v}$. For pure states $\ket{\psi}$ and $\ket{\varphi}$, the trace norm and the Euclidean norm satisfy the inequalities:
\begin{align}
\label{eq:relation of Euclidean and trace norm}
    \f{1}{2}\big\|\ketbra{\psi}{\psi} - \ketbra{\varphi}{\varphi}\big\|_1 
    \leq \min_\theta\big\|\ket{\psi} - e^{i\theta}\ket{\varphi}\big\| \leq \f{1}{\sqrt{2}}\big\|\ketbra{\psi}{\psi} - \ketbra{\varphi}{\varphi}\big\|_1.
\end{align}


%======================================================================================================================================



\subsection{Uhlmann transformation}
\label{sec:uhl intro}

We first recall the Uhlmann's theorem~\cite{uhlmann1976transition}.
\begin{theorem}[Uhlmann's theorem~\cite{uhlmann1976transition}]
\label{thm:Uhlmann theorem}
    Let $\rho^\sA$ and $\sigma^\sA$ be quantum states. The fidelity $\rF(\rho^\sA, \sigma^\sA)$ can be rephrased using their purified states $\ket{\rho}^{\sA\sB}$ and $\ket{\sigma}^{\sA\sB}$ as
\begin{align}
    \rF(\rho^\sA, \sigma^\sA) 
    &= \max_{U^\sB}\big| \bra{\sigma}^{\sA\sB}(\bI^\sA \otimes U^{\sB})\ket{\rho}^{\sA\sB} \big|^2 \\
    &= \max_{U^\sB}\rF(U^\sB \ket{\rho}^{\sA\sB}, \ket{\sigma}^{\sA\sB}),
\end{align}
where the maximization is taken over all unitaries acting on $\sB$.
\end{theorem}
This theorem states that for a given state $\ket{\rho}^{\sA\sB}$, there exists a unitary $U^\sB$ that realizes the transformation from $\ket{\rho}^{\sA\sB}$ to $\ket{\sigma}^{\sA\sB}$, with the minimum error achievable by local operations on $\sB$.
This optimal transformation and the corresponding unitary $U^\sB$ are called the \emph{Uhlmann transformation} and the \emph{Uhlmann unitary}, respectively.


While we focus here on transformations between pure states, the Uhlmann transformation actually characterizes the optimal local transformation between mixed states, by considering an environment. See \Cref{sec:opt local and Uhlmann}.


We should note several remarks. 
First, since we can always expand the purifying systems of $\rho^\sA$ and $\sigma^\sA$ by adding additional qubits, we can assume that their purifying systems are of the same size, which we denote by $\sB$.
Second, the Uhlmann unitary is not unique; in the kernel of $\rho^\sB$, the unitary $U^\sB$ can act arbitrarily.
Hence, the partial isometry $V^\sB$, which is uniquely determined as the part of $U^\sB$ that acts only on the support of $\rho^\sB$, is truly crucial.
We refer to this partial isometry $V^\sB$ as a \emph{Uhlmann partial isometry}.

An explicit form of the Uhlmann partial isometry $V^\sB$, which consists of $\ket{\rho}^{\sA\sB}$ and $\ket{\sigma}^{\sA\sB}$, is implied in Ref.~\cite{Jozsa1994FidelityforMixedQuantumStates}, and recently its uniqueness and robustness have been studied in Ref.~\cite{bostanci2025localtransformationsbipart}.
To present this, we define the sign function over real numbers as  
\begin{equation}
\sgn(x) = 
\begin{cases}
-1 & x<0, \\
0 & x=0, \\
1 & x>0.
\end{cases}
\end{equation}
Then, for a matrix $A$, $\sgn^{(\rm SV)}(A)$ is defined as the sign function acting on the singular values of $A$; that is, when $A$ has the singular value decomposition $A = \sum_j s_j \ketbra{\eta_j}{\xi_j}$, we have
\begin{align}
\label{inteq:158}
    \sgn^{(\rm SV)}(A) = \sum_j \sgn(s_j) \ketbra{\eta_j}{\xi_j}.
\end{align}



\begin{proposition}[Explicit form of the Uhlmann partial isometry~{\cite[Lemma 6]{Jozsa1994FidelityforMixedQuantumStates}}]
\label{prop:explicit form of Uhl}
An explicit form of the Uhlmann partial isometry is given by 
\begin{align}
\label{eq:explicit form Uhl partial iso}
    V^\sB = \sgn^{(\rm SV)}\big(\tr_\sA\big[\ketbra{\sigma}{\rho}^{\sA\sB}\big]\big).
\end{align}
\end{proposition}




One can readily check that the singular values of $\tr_\sA\big[\ketbra{\sigma}{\rho}^{\sA\sB}\big]$ and $\sqrt{\sigma^\sA}\sqrt{\rho^\sA}$ are identical.
Their ranks are also equal, as they have the same number of non-zero singular values.
In \Cref{sec:analyze partial iso}, we provide a derivation of \Cref{prop:explicit form of Uhl} for completeness.




Note that, since $V^\sB$ in Eq.~\eqref{eq:explicit form Uhl partial iso} is a partial isometry, it cannot be directly implemented as a deterministic physical operation. However, with the assistance of an ancillary system, the Uhlmann transformation can be approximately realized as a quantum channel.
That is, we aim for a transformation: $\ket{\rho}^{\sA\sB}\ket{\phi}^\sF \mapsto V^\sB\ket{\rho}^{\sA\sB}\ket{\phi'}^\sF$, with auxiliary states $\ket{\phi}^\sF$ and $\ket{\phi'}^\sF$, by operating on the system $\sB\sF$.






%=============================================================================================================================================================================





\subsection{Computational models for algorithm design}
\label{sec:three computational models}


We construct quantum algorithms that realize the Uhlmann transformation in the following three standard computational models. 
Here, we explain these computational models for a general state $\omega$; in the case of the Uhlmann transformation, $\omega$ is taken to be $\rho$ and $\sigma$.
\begin{enumerate}[label=\Roman*.]
    \item Purified query access model~\cite{watrous2002limitQSZK, watrous2009QSZKquantumattacks, belovs2019ClassicalDist, gilyen2019DistPropTesting, brandao2019SDPOpt, vanapeldoorn2019impleveSDP, Subramanian2021estimatealphalenni, gilyen2022improvedfidelity, Apeldoorn2023TomoPurifQuery, wang2023QAfidelityest, wang2023SampletoQuary, luo2024SuccinctTest, wang2024NewQuantEnt, Liu2024geometricmean}: let $U_\omega^{\sA\sB}$ be a unitary that prepares its purified state $\ket{\omega}^{\sA\sB}$, i.e., $U_\omega^{\sA\sB}\ket{0}^{\sA\sB}=\ket{\omega}^{\sA\sB}$.
    We assume that $U_\omega^{\sA\sB}$ is available multiple times as a unitary oracle.
    It is also assumed that we can utilize its inverse $(U_\omega^{\sA\sB})^\dag$.
    \item Purified sample access model~\cite{chen2024localtestUniInv, liu2024exponentialseparations}:
    multiple independent and identical copies of purified states $\ket{\omega}^{\sA\sB}$ are available.
    \item Mixed sample access model~\cite{gilyen2022improvedfidelity, wang2023SampletoQuary, wang2024TimeEfficientEntEstSamp, wang2024sampleoptimal}: multiple independent and identical copies of states $\omega^\sA$ are available.
\end{enumerate}




The purified query access model (Model~I) has been actively investigated in connection with quantum algorithms for processing classical-data distributions in oracle settings~\cite{belovs2019ClassicalDist, gilyen2019DistPropTesting, luo2024SuccinctTest}.
It is also commonly employed in quantum computational complexity~\cite{watrous2002limitQSZK, watrous2009QSZKquantumattacks}, quantum query algorithms~\cite{gilyen2022improvedfidelity, Apeldoorn2023TomoPurifQuery}, and semidefinite programming~\cite{brandao2019SDPOpt, vanapeldoorn2019impleveSDP}. 
The purified sample access model (Model~II) was recently considered in Refs.~\cite{chen2024localtestUniInv, liu2024exponentialseparations}.
Model~I and Model~II are suitable when we can manipulate the entire system, for example, in simulating quantum systems on a quantum computer.
Model~II further allows the relevant states to be prepared by other parties, even when one is unable to generate them. 
The mixed sample access model (Model~III) is now the standard sample input model for quantum property testing (see, e.g., Ref.~\cite{MdW16surveyquantproptest}), employed in many fundamental tasks such as quantum state tomography~\cite{ODonnell2016EffTomo, Haah2017samploptTomo, Guta2020FastTomo, Chen2023WhenAdapTomo, hu2024sampleoptimalmemoryefficient}. 


These computational models are hierarchically related in terms of their computational power.
It is clear that Model~III can be simulated by Model~II, and Model~II can be simulated by Model~I. 
In other words, let $\mathrm{Q}$, $\mathrm{S}_{\mathrm{p}}$, and $\mathrm{S}_{\mathrm{m}}$ be the query/sample complexity of Models I, II, and III, respectively. 
Then, it immediately holds that $\mathrm{Q} \leq \mathrm{S}_{\mathrm{p}} \leq \mathrm{S}_{\mathrm{m}}$ for any computational tasks. 
However, their quantitative relation is generally unclear. 
A known simple separation between $\mathrm{Q}$ and $\mathrm{S}_{\mathrm{p}}$ is by, for example, the estimation of pure-state closeness: for the problem of estimating the trace distance between two pure states to within additive error $\delta$, it is known that $\mathrm{Q} = \Theta(1/\delta)$~\cite{Wang2024OptimalFidelityPure,Fang2025OptimalFidelityToPure} whereas $\mathrm{S}_{\mathrm{p}} = \Omega(1/\delta^2)$ (implied by Ref.~\cite{ALL22} using quantum state discrimination). 
On the other hand, it was shown in Ref.~\cite{chen2024localtestUniInv} that, for testing unitarily invariant properties, $\mathrm{S_p} = \mathrm{S_m}$ holds, that is, Model~II has no advantage over Model~III.
In general cases, however, whether there is a strict separation between $\mathrm{S}_{\mathrm{p}}$ and $\mathrm{S}_{\mathrm{m}}$ is not yet known.

Our results provide some insights into this open problem.
As can be seen in \Cref{tab:compare Uhlmann and tomography}, our algorithms for the Uhlmann transformation for Models~II and~III come with distinct sample complexity.
Although this still does not prove the separation between $\mathrm{S}_{\mathrm{p}}$ and $\mathrm{S}_{\mathrm{m}}$ due to the lack of a good lower bound for $\mathrm{S}_{\mathrm{m}}$, our results suggest that the Uhlmann transformation is a viable candidate that could show the separation.






%=============================================================================================================================================================================





\subsection{Building blocks of quantum algorithms}
\label{sec:previous key subroutine}


The Uhlmann transformation algorithm can be constructed by combining several quantum algorithms as subroutines. 
Here, we briefly review these primitives: the block-encoding and the QSVT in \Cref{sec:BE and QSVT intro}, and the density matrix exponentiation in \Cref{sec:DME intro}.


\subsubsection{Block encoding and the quantum singular value transformation}
\label{sec:BE and QSVT intro}



The notion of block-encoding was introduced in Refs.~\cite{Low2019HamiltonianQubitize, gilyen2019qsvt} to describe that certain operators of interest $A$ are encoded as a block of a larger unitary $U$.
The block is specified by an $a$-qubit auxiliary state $\ket{0^a}$ as 
\begin{align}
        U
     = \hspace{1mm}
\begin{blockarray}{ccc}
& \bra{0^a} &  & \vspace{1mm}\\
\begin{block}{c(cc)}
  \ket{0^a} \hspace*{1mm} & A/\alpha & \hspace{0mm} * \hspace{3mm}\\
   \hspace*{1mm} & * & \hspace{0mm} * \hspace{3mm}\\
\end{block}
\end{blockarray}
\   \iff   A = \alpha(\bra{0^a} \otimes \bI)U(\ket{0^a} \otimes \bI).
\end{align}
Here, $\alpha$ is the normalization constant such that $\alpha \geq \|A\|_\infty$, which ensures that $U$ is a unitary with appropriate choices of the remaining blocks in $U$. 
More generally, a unitary $U$ is called an $(\alpha, a, \epsilon)$-block-encoding (unitary) of $A$, when $\big\|A - \alpha(\bra{0^a} \otimes \bI)U(\ket{0^a} \otimes \bI)\big\|_\infty \leq \epsilon$. If $(\alpha, a, \epsilon) = (1, a, 0)$ and $a$ is clear from the context, we simply refer to $U$ as an (exact) block-encoding (unitary) of $A$ and omit $a$ from $\ket{0^a}$.





Block-encoding is a powerful tool for implementing flexible operations.
In particular, one can construct a unitary $\tilde{U}$ that block-encodes a certain degree odd or even polynomial $P$ of a matrix.  
The construction is called the \emph{quantum singular value transformation (QSVT)}~\cite{gilyen2019qsvt, Gilyn2019thesis}:
\begin{align}
        U
     = \!\hspace{1mm}
\begin{blockarray}{ccc}
& \bra{0^a} &  & \vspace{1mm}\\
\begin{block}{c(cc)}
  \ket{0^a} \hspace*{1mm} & A/\alpha & \hspace{0mm} * \hspace{3mm}\\
   \hspace*{1mm} & * & \hspace{0mm} * \hspace{3mm}\\
\end{block}
\end{blockarray}
\ \ \overset{\rm QSVT}{\longrightarrow}  \
        \til{U}
     = \!\hspace{1mm}
\begin{blockarray}{ccc}
& \bra{0^{\til{a}}} &  & \vspace{1mm}\\
\begin{block}{c(cc)}
  \ket{0^{\til{a}}} \hspace*{1mm} & P^{(\rm SV)}(A/\alpha) & \hspace{0mm} * \hspace{3mm}\\
   \hspace*{1mm} & * & \hspace{0mm} * \hspace{3mm}\\
\end{block}
\end{blockarray}\!\hspace{2.5mm},
\end{align}
where $P^{(\rm SV)}(A/\alpha)$ is defined by 
\begin{align}
\label{eq:P poly def}
    P^{(\rm SV)}(A/\alpha) =
    \begin{cases}
    \sum_j P(s_j/\alpha)\ketbra{\eta_j}{\xi_j}, \ \  \text{if $P$ is odd,} \\
    \sum_j P(s_j/\alpha)\ketbra{\xi_j}{\xi_j}, \ \ \text{if $P$ is even,}
    \end{cases}
\end{align}
and $A = \sum_j s_j \ketbra{\eta_j}{\xi_j}$ is the singular value decomposition of $A$.
The formal statement of the QSVT, specialized to real even/odd polynomials, is as follows, where the polynomial $P$ must satisfy some conditions.


\begin{theorem}[Quantum singular value transformation with a real polynomial~{\cite[Corollary 11]{gilyen2019qsvt}}; see also Ref.~{\cite[Corollary 2.3.8]{Gilyn2019thesis}}]
\label{prop:general QSVT}
Let $P$ be a real even/odd polynomial of degree-$l$ satisfying $|P(x)| \leq 1$ for $x \in [-1,1]$, and let $U$ be an $(\alpha, a, 0)$-block-encoding unitary of $A$.
Then, a quantum circuit for implementing a block-encoding unitary $\tilde{U}$ of $P^{(\rm SV)}(A/\alpha)$ can be constructed using $\f{l+1}{2}$ calls to $U$ and $\f{l-1}{2}$ calls to $U^\dag$ if $P$ is odd, and $\f{l}{2}$ calls to $U$ and $\f{l}{2}$ calls to $U^\dag$ if $P$ is even, along with $\cO(al)$ other one- and two-qubit gates.  
\end{theorem}

Note that it is straightforward to construct $\tilde{U}^\dag$ by slightly modifying the construction of $\tilde{U}$~\cite{gilyen2019qsvt, Gilyn2019thesis, martyn2021grand}.


By choosing a polynomial $P$ appropriately, we can implement various useful quantum algorithms within the QSVT framework, such as the amplitude amplification, matrix inversion, and Hamiltonian simulation.
To implement the QSVT in practice, one needs to compute a sequence of real phase factors corresponding to a given $P$. Algorithms that compute such a phase factor sequence in time polynomial in the degree of $P$ are known~\cite{Haah2019product, Chao2020FindingAF, dong2021findingphase, Lin2022LectureNotes, Mizuta2023vzw}.







The following proposition is a result of selecting an appropriate real odd polynomial $P_\sgn$ that closely approximates the sign function~\cite{gilyen2019qsvt, Gilyn2019thesis, martyn2021grand}.
For future reference, we present it in an algorithmic form.

\begin{proposition}[QSVT with the sign function~{\cite[Theorem 1]{gilyen2019qsvt}}; see also Ref.~{\cite[Theorem 3.2.3]{Gilyn2019thesis}}]
\label{prop:FPAA general}
    Let $\delta \in (0, 1/2)$ and $\beta \in (0, 1]$, and let $U$ be a $(1, a, 0)$-block-encoding unitary of a matrix $M$. Then, there exists a quantum algorithm $\mathtt{QSVTSIGN}(U, U^\dag; \delta, \beta)$ that outputs a unitary $\til{U}$, which is a $(1, a+1, 0)$-block-encoding of a degree-$u$ polynomial $P_\sgn^{(\rm SV)}(M)$ such that $\big|P_\sgn(x) - \sgn(x)\big| \leq \delta$,
    for $|x| \in [\beta, 1]$ and $|P_\sgn(x)| \leq 1$ for $|x| \in [0, 1]$, where $u$ is the minimum odd integer satisfying $u \geq \big\lceil\f{8e}{\beta}\log{(2/\delta)}\big\rceil$.
    This algorithm uses $U$ and $U^\dag$ a total of $u$ times.
\end{proposition}


This proposition is based on the fact that, for given $\delta$ and $\beta$, there exists a degree-$u$ polynomial $P_{\rm sgn}$ that well approximates the sign function.
More precisely, for any odd degree $u \geq \gamma_{\delta, \beta}$, where
\begin{align}
    &\gamma_{\delta, \beta} = 2 \Big\lceil\max\Big\{
    \frac{e}{2\beta} \sqrt{W\Big(\frac{8}{\pi\delta^2}\Big) W\Big(\frac{512}{e^2\pi\delta^2}\Big)}, \sqrt{2}\,W\Big(\frac{4}{\beta\delta}\sqrt{\frac{2}{\pi}W\Big(\frac{8}{\pi\delta^2}\Big)}\Big)\Big\}\Big\rceil + 1,
\end{align}
there is a polynomial that approximates $\sgn(x)$ within error $\delta$ over $|x| \in [\beta, 1]$~\cite{martyn2023efficient}.
Here, $W(x)$ is the Lambert $W$ function~\cite{Corless1996onthelambert}.
From the fact that $W(x) \leq \log x$ holds for $x > e$~\cite{Hassani2005Approximationlambert, Hoorfar2008inequlambert}, we have $\gamma_{\delta, \beta} \leq \big\lceil\frac{8e}{\beta}\log(2/\delta)\big\rceil$ for $\delta \in (0,1/2)$ and $\beta \in (0,1]$. 
Thus, it is sufficient to set the number of uses $u$ of the unitaries $U$ and $U^\dag$, which coincides with the degree of the polynomial, to the minimum odd integer satisfying $u \geq \big\lceil\frac{8e}{\beta}\log(2/\delta)\big\rceil$.











\subsubsection{Density matrix exponentiation}
\label{sec:DME intro}



Given identical copies of a quantum state $\rho$, we can approximately implement a unitary $e^{it\rho}$ using a technique known as the \emph{density matrix exponentiation}, which was introduced in Ref.~\cite{lloyd2014QuantPrincCompAnaly} and developed in subsequent works~\cite{marvian2016universalquantumemulator, kimmel2017HamSimSampleComp, wei2023DMEhermitianpreserv, go2024DMEsamplebased}.


\begin{theorem}[Density matrix exponentiation~{\cite[Theorem 2]{go2024DMEsamplebased}}]
\label{prevthm:cntrl-DME}
Let $\rho$ be a quantum state in a $d$-dimensional Hilbert space.
For $\delta > 0$ and $t \geq \delta/4$, there exists a quantum algorithm that realizes a quantum channel $\cG$ such that
$\f{1}{2}\|\cG - \cU\|_\diamond \leq \delta$,
where $\cU(\cdot) = e^{it\rho}(\cdot)e^{-it\rho}$.
The algorithm uses $m = \lceil4t^2/\delta\rceil$ copies of $\rho$, and $\cO(m\log{d})$ one- and two-qubit gates.
\end{theorem}



The key idea of the density matrix exponentiation is as follows.
If, for a small parameter $\Delta t$, we apply partial swap unitary $e^{i\Delta t F^{\sA\hat{\sA}}}$, we see that
\begin{align}
    \tr_{\hat{\sA}}\big[e^{i\Delta t F^{\sA\hat{\sA}}} (\psi^{\sR\sA} \otimes \rho^{\hat{\sA}}) e^{-i\Delta t F^{\sA\hat{\sA}}}\big] 
    &= \psi^{\sR\sA} + i\Delta t [\rho^\sA, \psi^{\sR\sA}] + \cO\big((\Delta t)^2\big) \\ 
    &= e^{i\Delta t\rho^\sA} \psi^{\sR\sA} e^{-i\Delta t \rho^\sA} + \cO\big((\Delta t)^2\big),
\end{align}
where the last line follows from the Hadamard lemma, also known as the Baker-Campbell-Hausdorff formula~\cite{campbell1896BCHlemma, Baker1901BCHlemma, hausdorff1906symbolische, dynkin1947calculation}.
By repeating this procedure $m$ times, we end up implementing the unitary $e^{i\rho^\sA t}$ up to the error $\cO\big(m(\Delta t)^2)$, where $m \Delta t = t$. 
Thus, to achieve the approximation of $e^{it\rho^\sA}$ within an error of $\delta$, it is sufficient to set $m = \Omega(t^2/\delta)$, which represents the number of copies of $\rho^\sA$.
A quantum channel that approximates $\cU^\dag(\cdot) = e^{-it\rho}(\cdot)e^{it\rho}$ can also be constructed by simply modifying this procedure to replace $e^{i\Delta t F}$ with $e^{-i\Delta t F}$.


Intuitively, the first-order relation: $\tr_{\hat{\sA}}\big[F^{\sA\hat{\sA}}(\psi^{\sR\sA}\otimes\rho^{\hat{\sA}})\big] = \rho^\sA\psi^{\sR\sA}$, implies that by repeatedly preparing $\rho^{\hat{\sA}}$ and applying $e^{i\Delta t F^{\sA\hat{\sA}}}$ and $\tr_{\hat{\sA}}$, higher-order deviations $\delta$ can be sufficiently suppressed, which enables us to approximately simulate the unitary $e^{it\rho^\sA}$.
For detailed discussions of higher-order terms, refer to Refs.~\cite{wei2023DMEhermitianpreserv, go2024DMEsamplebased}.



The density matrix exponentiation approximately realizes a variety of unitary operations by appropriately preparing multiple copies of states and designing repeated operations.
The following is one of the results, provided in an algorithmic form for future convenience.


\begin{proposition}[Density matrix exponentiation for the difference of two subnormalized states~{\cite[Lemma 12]{kimmel2017HamSimSampleComp}}]
\label{prevthm:DME variant}
Let $\delta > 0$ and $t \geq \delta/4$, and let $\Upsilon$ be a quantum state in a $d$-dimensional Hilbert space of the form  
\begin{align}
    \Upsilon = \ketbra{0}{0} \otimes \xi_0 + \ketbra{1}{1} \otimes \xi_1,
\end{align}
where $\xi_0$ and $\xi_1$ are subnormalized states with $\tr[\xi_0] + \tr[\xi_1] = 1$. 
Then, there exists a quantum algorithm $\mathtt{DMESUB}(\Upsilon;\delta, t)$ that outputs a quantum channel $\cG$ such that $\frac{1}{2} \big\|\cG - \cU\big\|_\diamond \leq \delta$, where $\cU(\cdot) = e^{it(\xi_0 - \xi_1)}(\cdot)e^{-it(\xi_0 - \xi_1)}$.
The algorithm uses $m = \lceil 4t^2/\delta\rceil$ samples of $\Upsilon$ and $\cO(m\log{d})$ one- and two-qubits gates.

Moreover, there exists a quantum algorithm $\mathtt{DMESUB}^\dag(\Upsilon; \delta, t)$ that outputs a quantum channel $\cG_{\rm inv}$ satisfying $\frac{1}{2} \big\|\cG_{\rm inv} - \cU^\dag \big\|_\diamond \leq \delta$, using the same number of samples and gates as $\mathtt{DMESUB}(\Upsilon; \delta, t)$.
\end{proposition}


The main difference from \Cref{prevthm:cntrl-DME} is that, instead of using the partial swap $e^{i\Delta t F}$, this algorithm adopts a unitary operator $\ketbra{0}{0} \otimes e^{i\Delta t F} + \ketbra{1}{1} \otimes e^{-i\Delta t F}$.
When we see the first-order term, we find that
\begin{align}
    \tr_{\sC\hat{\sA}}\big[\big(\ketbra{0}{0}^\sC \otimes F^{\sA\hat{\sA}}
    + \ketbra{1}{1}^\sC \otimes (-F^{\sA\hat{\sA}})\big) \big(\psi^{\sR\sA}\otimes (\ketbra{0}{0}^\sC \otimes \xi_0^{\hat{\sA}} + \ketbra{1}{1}^\sC \otimes \xi_1^{\hat{\sA}})\big)\big]
    = (\xi_0^\sA - \xi_1^\sA)\psi^{\sR\sA}.
\end{align}
From its similarity to the standard density matrix exponentiation, we can intuitively understand that the operation realizes $e^{it(\xi_0^\sA - \xi_1^\sA)}$.
By only replacing $e^{i\Delta t F}$ with $e^{-i\Delta t F}$, we can implement the inverse $e^{-it(\xi_0^\sA - \xi_1^\sA)}$ using the rest of the same procedure.






%==================================================================================================================================================
